\chapter{Methodology}
\label{ch:methodology}

% ============================================================================
% This chapter describes the methodology: dataset, preprocessing, features, models
% ============================================================================

\section{Overview}
\label{sec:method_overview}

This chapter describes the methodology for sleep stage classification. Figure~\ref{fig:pipeline_overview} shows the overall pipeline:

\begin{enumerate}
    \item \textbf{Data Loading:} Load raw EEG data from BOAS dataset
    \item \textbf{Preprocessing:} Filter, artifact removal, epoch segmentation
    \item \textbf{Feature Extraction:} Extract 149 features per epoch
    \item \textbf{Feature Selection:} ANOVA-based selection with correlation filtering
    \item \textbf{Model Training:} Train XGBoost/Random Forest with LOSO cross-validation
    \item \textbf{Evaluation:} Compute accuracy, kappa, F1-score
\end{enumerate}

% TODO: Add pipeline diagram (Figure 1)
% \begin{figure}[htbp]
%     \centering
%     \includegraphics[width=\textwidth]{figures/pipeline_overview.pdf}
%     \caption{Sleep stage classification pipeline overview}
%     \label{fig:pipeline_overview}
% \end{figure}

\section{Dataset: BOAS}
\label{sec:dataset}

\subsection{Dataset Description}
\label{subsec:dataset_description}

The BOAS dataset contains polysomnography recordings from 128 healthy subjects. Each subject has one full night of sleep recorded with 6 EEG channels. Table~\ref{tab:boas_summary} summarizes the dataset.

% TODO: Add table with dataset statistics
\begin{table}[htbp]
    \centering
    \caption{BOAS Dataset Summary}
    \label{tab:boas_summary}
    \begin{tabular}{ll}
        \toprule
        \textbf{Property} & \textbf{Value} \\
        \midrule
        Number of subjects & 128 \\
        EEG channels & 6 \\
        Sampling rate & 100 Hz \\
        Epoch duration & 30 seconds \\
        Total epochs & $\sim$119,763 \\
        Average sleep duration & $\sim$7.8 hours/subject \\
        Sleep stages & 5 (W, N1, N2, N3, REM) \\
        \bottomrule
    \end{tabular}
\end{table}

\subsection{EEG Channels}
\label{subsec:eeg_channels}

The 6 EEG channels are:
\begin{itemize}
    \item C3-A2
    \item C4-A1
    \item F3-A2
    \item F4-A1
    \item O1-A2
    \item O2-A1
\end{itemize}

These channels cover frontal, central, and occipital brain regions, providing comprehensive spatial coverage for sleep stage classification.

\subsection{Sleep Stage Distribution}
\label{subsec:stage_distribution}

% TODO: Add bar chart of sleep stage distribution
% Figure showing class imbalance (N2 is typically most common)

The dataset exhibits class imbalance typical of sleep data, with N2 being the most common stage and N1 being the rarest.

\section{Preprocessing}
\label{sec:preprocessing}

\subsection{Signal Filtering}
\label{subsec:filtering}

We apply a bandpass filter from 0.5 Hz to 30 Hz to remove low-frequency drift and high-frequency noise. This range captures the relevant EEG frequencies for sleep staging:

\begin{itemize}
    \item \textbf{Delta (0.5--4 Hz):} Dominant in deep sleep (N3)
    \item \textbf{Theta (4--8 Hz):} Prominent in N1 and REM
    \item \textbf{Alpha (8--13 Hz):} Appears during drowsiness and wake
    \item \textbf{Beta (13--30 Hz):} Active during wake state
\end{itemize}

\subsection{Epoch Segmentation}
\label{subsec:segmentation}

The continuous EEG signal is segmented into 30-second epochs following standard sleep scoring practice \cite{TODO_ADD_AASM}. Each epoch is labeled with the corresponding sleep stage.

\subsection{Artifact Handling}
\label{subsec:artifacts}

% TODO: Describe if any artifact rejection is performed
% Check preprocessing.py for details

\section{Feature Extraction}
\label{sec:feature_extraction}

We extract 149 features per epoch, combining time-domain and frequency-domain characteristics from all 6 EEG channels.

\subsection{Time-Domain Features}
\label{subsec:time_features}

Time-domain features capture statistical properties of the signal:

\begin{itemize}
    \item \textbf{Mean:} Average amplitude
    \item \textbf{Variance:} Signal variability
    \item \textbf{Skewness:} Distribution asymmetry
    \item \textbf{Kurtosis:} Distribution tail heaviness
    \item \textbf{Zero-crossing rate:} Frequency indicator
    \item \textbf{Line length:} Signal complexity measure
    \item \textbf{Hjorth parameters:} Activity, mobility, complexity
\end{itemize}

\subsection{Frequency-Domain Features}
\label{subsec:freq_features}

Frequency-domain features are extracted using Welch's method for power spectral density estimation:

\begin{itemize}
    \item \textbf{Band power:} Absolute power in delta, theta, alpha, beta bands
    \item \textbf{Relative band power:} Power normalized by total power
    \item \textbf{Band power ratios:} Delta/alpha, theta/beta, etc.
    \item \textbf{Spectral entropy:} Frequency distribution uniformity
    \item \textbf{Spectral edge frequency:} Frequency below which X\% of power resides
\end{itemize}

\subsection{Multi-Channel Features}
\label{subsec:multichannel}

For each feature type, we compute values across all 6 channels, resulting in 149 total features per epoch.

% TODO: Add feature count breakdown table

\section{Feature Selection}
\label{sec:feature_selection_method}

Feature selection reduces dimensionality and improves model generalization. We employ a two-stage approach:

\subsection{Correlation Filtering}
\label{subsec:correlation_filtering}

First, we remove highly correlated features ($r > \rho$, where $\rho \in \{0.90, 0.95, 0.99\}$) to reduce redundancy. For each pair of correlated features, we retain the one with higher ANOVA F-score.

\subsection{ANOVA F-Test Selection}
\label{subsec:anova}

Second, we rank features using ANOVA F-test and select the top-$k$ features, where $k \in \{30, 50, 75\}$ in our experiments.

ANOVA F-test measures the ratio of between-class variance to within-class variance:

\begin{equation}
F = \frac{\text{Between-class variance}}{\text{Within-class variance}}
\end{equation}

Higher F-scores indicate features with better class separation.

\subsection{Feature Selection Scope}
\label{subsec:selection_scope}

% TODO: Clarify if feature selection is global or per-fold
% Check feature_selection.py for current implementation

\section{Machine Learning Models}
\label{sec:models}

We evaluate two tree-based ensemble models:

\subsection{XGBoost}
\label{subsec:xgboost}

XGBoost (Extreme Gradient Boosting) \cite{TODO_ADD_XGBOOST_PAPER} is a gradient boosting framework that builds an ensemble of decision trees sequentially, with each tree correcting errors of the previous ones.

Key hyperparameters:
\begin{itemize}
    \item \textbf{max\_depth:} [value]
    \item \textbf{learning\_rate:} [value]
    \item \textbf{n\_estimators:} [value]
    \item \textbf{subsample:} [value]
\end{itemize}

% TODO: Fill in hyperparameter values from config.py

\subsection{Random Forest}
\label{subsec:random_forest}

Random Forest \cite{TODO_ADD_RF_PAPER} builds an ensemble of decision trees trained on random subsets of features and samples, combining predictions through majority voting.

Key hyperparameters:
\begin{itemize}
    \item \textbf{n\_estimators:} [value]
    \item \textbf{max\_depth:} [value]
    \item \textbf{min\_samples\_split:} [value]
\end{itemize}

% TODO: Fill in hyperparameter values from config.py

\section{Leave-One-Subject-Out Cross-Validation}
\label{sec:loso_method}

We use Leave-One-Subject-Out (LOSO) cross-validation to evaluate subject-independent performance. For each of the 128 subjects:

\begin{enumerate}
    \item Hold out one subject as test set
    \item Train model on remaining 127 subjects
    \item Evaluate on held-out subject
    \item Repeat for all 128 subjects
\end{enumerate}

Final metrics are averaged across all 128 folds. This ensures models generalize to unseen individuals, simulating real-world deployment.

\section{Evaluation Metrics}
\label{sec:metrics}

We evaluate models using:

\begin{itemize}
    \item \textbf{Accuracy:} Overall percentage of correct predictions
    \item \textbf{Cohen's Kappa ($\kappa$):} Chance-corrected agreement
    \item \textbf{F1-Score:} Harmonic mean of precision and recall (macro-averaged across classes)
    \item \textbf{Per-Class Metrics:} Precision, recall, F1 for each sleep stage
    \item \textbf{Confusion Matrix:} Detailed breakdown of predictions
\end{itemize}

\section{Experimental Design}
\label{sec:experimental_design}

Our experimental grid consists of:

\begin{itemize}
    \item \textbf{Models:} XGBoost, Random Forest (2 models)
    \item \textbf{Correlation thresholds:} 0.90, 0.95, 0.99 (3 values)
    \item \textbf{Top-K features:} 30, 50, 75 (3 values)
\end{itemize}

This yields $2 \times 3 \times 3 = 18$ configurations. With LOSO cross-validation on 128 subjects, this requires $18 \times 128 = 2{,}304$ model training runs.

% TODO: Clarify if thesis uses simplified grid (2 × 3 = 6 configs)
% Based on earlier context, thesis uses 2 models × 3 top-K = 6 configs

\section{Summary}
\label{sec:method_summary}

This chapter described the methodology: BOAS dataset with 128 subjects and 6 EEG channels, preprocessing with bandpass filtering and epoch segmentation, extraction of 149 features, ANOVA-based feature selection, and training of XGBoost and Random Forest models with LOSO cross-validation. The next chapter details the implementation.
